% Standalone note: compile this file directly with pdfLaTeX or latexmk.
% No manuscript inputs or files from the DIGRAM repository are required.
\documentclass[11pt,a4paper]{article}
\usepackage[T1]{fontenc}
\usepackage[utf8]{inputenc}
\usepackage[english]{babel}
\usepackage{lmodern,microtype}
\usepackage[margin=25mm]{geometry}
\usepackage{amsmath,amssymb,booktabs,tabularx,longtable,array}
\PassOptionsToPackage{obeyspaces}{url}
\usepackage{enumitem,needspace,xcolor,xurl}
\usepackage[hidelinks]{hyperref}
\hypersetup{pdftitle={Manuscript screening procedure and DIGRAM SCREEN J},
  pdfsubject={Detailed source-based algorithm comparison}}
\definecolor{tracegray}{gray}{0.32}
\setlength{\parskip}{4pt}
\setlength{\parindent}{0pt}
\setlength{\emergencystretch}{1.5em}
\setlist{itemsep=3pt,topsep=4pt}
\renewcommand{\arraystretch}{1.15}
\newcolumntype{L}[1]{>{\raggedright\arraybackslash}p{#1}}
\newcommand{\SRC}{\operatorname{SOURCE}}
\newcommand{\DIF}{\operatorname{DIF}}
\newcommand{\WPG}{\operatorname{WPG}}
\newcommand{\ind}{\perp}
\newcommand{\Iset}{\mathcal{I}}
\newcommand{\src}[2]{\hyperlink{source-#1}{\textsf{#1}}, lines~#2}
\newcommand{\trace}[2]{%
  \par\smallskip
  {\footnotesize\color{tracegray}\textbf{Manuscript:} #1.\quad
  \textbf{DIGRAM:} #2.\par}\medskip}
\newcommand{\finding}[1]{\Needspace{5\baselineskip}\subsection{#1}}
\newcommand{\man}{\textbf{Manuscript.}\ }
\newcommand{\dig}{\textbf{DIGRAM.}\ }
\newcommand{\effect}{\textbf{Consequence.}\ }

\title{Manuscript screening procedure and DIGRAM SCREEN J:\\
  a detailed algorithm comparison}
\author{Technical note for the coauthors}
\date{12 September 2026}

\begin{document}
\maketitle

\begin{abstract}
This note compares the six-step procedure in \emph{Item screening in graphical
Rasch models: a tutorial}, including its algorithm and data-example appendices,
with the original Pascal implementation of DIGRAM SCREEN J. The procedures
differ in local-dependence priorities, the updating of candidate sets,
differential-item-functioning purification, test statistics, score-effect
selection, and the scope of subsequent graph checks. The comparison uses the
manuscript's notation and gives source locations, explicit decision rules, and
examples of their consequences. DIGRAM's separate CHECK I routines are examined
where relevant to the manuscript's Step 6. Confirmed differences, missing
specifications, and internal example inconsistencies are distinguished.
\end{abstract}

\section{Scope, versions, and interpretation}
\label{sec:scope}
The manuscript baseline is PREP revision \texttt{330511f}
(\path{330511f4f1539c4dd971bbc51d6eea3581975ebe}).
The DIGRAM repository baseline is revision \texttt{cbc484a}
(\path{cbc484a43dae1b2356f851c7d4bc035f774e86c2}).
The original implementation examined is the preserved DIGRAM 7.04 snapshot
under \path{../digram/source/digram_source_20260817/}.
The project file \path{scd/scd.dpr} determines which duplicate unit is active:
the command handler comes from \path{scd/}, and the statistical units cited
here come from \path{skunits/}. Source abbreviations and paths are listed in
Section~\ref{sec:sources}.

The comparison concerns the behavior of this particular implementation.
A discrepancy can be an intentional methodological change; it does not,
by itself, establish that the manuscript's method is statistically inferior.
The R package used by the manuscript and the R reimplementation in the sister
repository are not substitutes for the original Pascal source in this audit.

Three types of finding appear below:
\begin{description}[style=nextline,leftmargin=0pt,labelsep=0pt,labelwidth=0pt]
\item[Behavioral difference.]
The two descriptions imply different tests, candidate sets, selection orders,
or outputs.
\item[Specification gap.]
The manuscript does not provide enough detail to determine whether its
behavior agrees with DIGRAM.
\item[Internal manuscript issue.]
A displayed example or pseudocode statement is inconsistent with another
part of the manuscript, independently of DIGRAM.
\end{description}

The numerical examples introduced specifically for this note are hypothetical
screening summaries, not results of newly fitted datasets. Statements about
the Actinic Keratosis Quality of Life (AKQoL) and Shoulder Pain and Disability
Index (SPADI) analyses identify consequences of the stated hypotheses
and graph structures. A complete numerical comparison would require rerunning
the analyses with matched data preparation and inference settings.

\subsection{The command boundary matters}
SCREEN J saves the current local-dependence (LD) and differential-item-functioning
(DIF) matrices, computes screening evidence, reports it, and restores the saved
matrices. The command handler encloses graph initialization, model definition,
and moralization in a branch that explicitly excludes command option J.
The boolean \texttt{DefineModel} passed through the screening wrapper is marked
as unused in the statistical routine; its name does not imply that SCREEN J
performs a final model check.

Consequently, manuscript Steps 5 and 6 are additional stages relative to
SCREEN J. Related graph diagnostics exist as separate DIGRAM commands,
including CHECK I and CHECK D. Those routines are compared explicitly in
Section~\ref{sec:step6}; their behavior must not be attributed to SCREEN J.
\trace{\path{article.tex}, lines 412--428 and 669--711}
{\src{CMD}{1793--1804, 1877--1922}; \src{EXE}{1351--1366};
 \src{S7}{2330--2339, 3228--3315}}

\section{Notation and the sequence being compared}
\label{sec:notation}
Let
\[
 \Iset=\{Y_1,\ldots,Y_k\},\qquad
 \boldsymbol X=(X_1,\ldots,X_m),\qquad
 S=\sum_{i=1}^kY_i,\qquad R_i=S-Y_i .
\]
For an item subset \(B\subseteq\Iset\), write
\[
 S_B=\sum_{Y_i\in B}Y_i,\qquad B^c=\Iset\setminus B .
\]
The \(Y_i\) are items and the \(X_j\) are exogenous covariates throughout this
note. Letters \(A,B,C,D,E\) and \(X,Z\) are retained when referring to the
manuscript's hypothetical examples. Write \(S_{\max}\) for the largest possible
total score; this avoids confusing it with the manuscript's per-item category
maximum \(M\) or its number of covariates \(m\).

At a given stage of screening, use the manuscript's set notation
\begin{align*}
 \SRC(Y_i)&=\{X_j:(Y_i,X_j)\text{ is currently retained}\},\\
 \DIF(X_j)&=\{Y_i:(Y_i,X_j)\text{ is currently retained}\}.
\end{align*}
To describe the graph checks without using \(\DIF\) for two different kinds
of object, extend the source notation explicitly:
\[
 \SRC(B)=\bigcup_{Y_i\in B}\SRC(Y_i),\qquad
 \mathcal B_{\mathrm{all}}=\bigcup_{j=1}^m\DIF(X_j).
\]
Thus \(\SRC(B)\) is a set of covariates, and
\(\mathcal B_{\mathrm{all}}\) is the set of items having any retained DIF.
These definitions clarify the comparison; they do not alter the manuscript.

For an unordered item pair \(\{Y_i,Y_{i'}\}\), its two directional tests are
\[
 H_{ii'}^{(i)}:Y_i\ind Y_{i'}\mid R_i,\qquad
 H_{ii'}^{(i')}:Y_i\ind Y_{i'}\mid R_{i'}.
\]
Write their estimates and p-values as \(\gamma_{ii'}^{(h)}\) and
\(p_{ii'}^{(h)}\), where \(h\in\{i,i'\}\). A direction identifies the
conditioning rest-score; it does not make statistical association asymmetric
within a fixed conditioning set.

Let \(c_{.05}\), \(c_{.01}\), and \(c_{.001}\) denote DIGRAM's raw-p-value
cutoffs obtained from its combined Benjamini--Hochberg (BH) family. These are
distinct from the nominal false discovery rate (FDR) levels \(0.05\),
\(0.01\), and \(0.001\), and from adjusted p-values.

\begin{center}
\begin{tabularx}{\linewidth}{@{}L{0.24\linewidth}X@{}}
\toprule
Manuscript stage & Corresponding original-source behavior\\
\midrule
Step 1 & Marginal diagnostics are reported; the manuscript's explicit
M1--M3 verdicts are an additional interpretive layer.\\
Step 2 & Two directed LD tests per item pair and one selected test statistic
per item--covariate DIF relation; one combined BH family.\\
Step 3(a) & Four successive LD-selection stages with sign-specific priorities
and original evidence counts.\\
Steps 3(b), 3(c) & DIGRAM executes the per-covariate item pass first
(manuscript 3(c)), then the per-item source pass (3(b)), using shared updates.\\
Step 4 & Marginal score-effect selection followed by conditional elimination.\\
Step 5 & SCREEN J reports evidence and restores the current model;
it does not construct the manuscript's new graphs.\\
Step 6 & No corresponding SCREEN J stage. Related CHECK I routines have
different conditioning and reporting behavior.\\
\bottomrule
\end{tabularx}
\end{center}

\section{Preparation and Step 1: consistency}
\label{sec:step1}
\finding{The operational scope is narrower than the measurement requirements}
\man M1 concerns positive associations between item pairs. M2 concerns
item--rest-score associations and all subscores excluding the item. M3 concerns
consistent directions between a covariate, the score, subscores, and items.
The Discussion explains that the implementation does not test all subscores.
The data examples give explicit pass/fail judgments and reasons to proceed.

\dig SCREEN J computes and prints marginal item--item, item--rest-score,
item--covariate, and score--covariate diagnostics. It does not enumerate all
item subsets. Its active orchestration does not stop or remove an item because
an explicit M1, M2, or M3 verdict fails. The local helper functions
\texttt{problem1} and \texttt{problem2} in \texttt{Item\_Screening} are
defined but are not called by that routine's active screening sequence;
their definitions should not be mistaken for executed pass/fail rules.

\effect The theoretical requirements, the manuscript's reduced checks,
and SCREEN J's diagnostics are three distinct levels of description.
Restricting checks to rest-scores is compatible with the source's scope,
but the manuscript's rules for interpreting failures and deciding to proceed
are not an automatic SCREEN J selection stage.
\trace{\path{article.tex}, lines 237--249, 565, and 702;
 \path{Appendix - data example.tex}, lines 13--15}
{\src{S7}{668--754, 2361--2378, 2652--2717, 2943--3008}}

\finding{Coding, record selection, and score restrictions need specification}
\man The examples use complete item and covariate data after preprocessing;
the SPADI appendix describes category collapsing. Algorithm 1 is stated for
a supplied dataset without a general rule for record frequencies, missing
covariates, score restrictions, or item reversal.

\dig SCREEN J requires at least three items. Source item codes are stored
with an offset, with configured reversals applied during item extraction;
the corresponding mathematical scores can be represented by the manuscript's
zero-based \(Y_i\). Frequency counts in DIGRAM tables contribute to all
contingency-table cells. The source distinguishes \texttt{GET\_ITEMS} from
\texttt{GET\_GOOD\_ITEMS}. The latter additionally requires
\[
 0<S<S_{\max},\qquad
 \texttt{MinScore}\le S\le\texttt{MaxScore}.
\]
The configured \texttt{MaxScore} is not silently identified with \(S_{\max}\).

For each initial LD pair \(i<i'\), the first direction uses
\texttt{GET\_ITEMS}; the reverse uses \texttt{GET\_GOOD\_ITEMS}.
Initial DIF tables use stored scores \(1,\ldots,\texttt{MaxScore}-1\),
original item/covariate codes, and a post-computation gamma sign reversal
for flipped items. Later DIF pruning rebuilds tables for the current
candidate set using valid non-extreme total scores.
Score-effect screening uses \texttt{GET\_ITEMS} and permits extreme scores.
The exact record set is therefore part of the computation.

\effect This is primarily a reproducibility gap. With ordinary complete data
and the full score range, some exclusions involve strata with no comparable
pairs and may leave gamma unchanged. That does not justify assuming all
source filters are immaterial under missingness or restricted score ranges.
Further table-building details are given in Section~\ref{sec:inference}.
\trace{\path{article.tex}, lines 140, 151, and 562;
 \path{Appendix - algorithms.tex}, lines 13--37;
 \path{Appendix - data example.tex}, lines 5--11}
{\src{CMD}{1779--1784}; \src{S2}{3318--3403};
 \src{S7}{267--314, 462--463, 777--815, 844--875}}

\section{Step 2: initial LD and DIF evidence}
\label{sec:step2}
\finding{The initial conditional hypotheses and pooled gamma agree}
The manuscript's initial hypotheses
\[
 Y_i\ind Y_{i'}\mid R_i,\qquad
 Y_i\ind Y_{i'}\mid R_{i'},\qquad
 Y_i\ind X_j\mid S
\]
agree with the source's conditioning structure. The main differences at this
stage concern the statistic, inference conventions, and multiplicity family.

The pooled partial gamma also agrees when its denominator is positive.
Writing \(C_\ell,D_\ell\) for concordant and discordant pair counts in
conditioning stratum \(\ell\),
\[
 \widehat\gamma=
 \frac{\sum_\ell(C_\ell-D_\ell)}
      {\sum_\ell(C_\ell+D_\ell)}.
\]
DIGRAM's \texttt{PMQ} and \texttt{PPQ} use doubled pair counts. The common
factor cancels in gamma and in the weighted directional gamma below.
\trace{\path{article.tex}, lines 253--259;
 \path{Appendix - association measures.tex}, lines 1--15}
{\src{S7}{759--875, 448--477}; \src{STAT}{1004--1030}}

\finding{Initial DIF is not always tested with gamma}
\man The Methods and algorithm appendix describe gamma and partial-gamma
tests for the screening procedure.

\dig Initial LD decisions use partial gamma. For initial DIF, the
covariate determines the selected statistic: binary covariates are treated
as ordinal, ordinal covariates use gamma, and other nominal covariates use
the source's contingency-table chi-square statistic.
The latter need not be Pearson's statistic: DIGRAM's global settings can
select a likelihood-ratio or other implemented power-divergence form.

\effect Gamma-only screening is a restricted or modified procedure.
For nominal covariates with more than two categories it can depend on an
arbitrary ordering that DIGRAM's initial nominal-variable branch does not require.
The example covariates avoid this particular nominal-multicategory issue,
but the manuscript's general algorithm does not state the restriction.
Later source stages treat variable types differently
(Section~\ref{sec:inference}).
\trace{\path{article.tex}, line 235;
 \path{Appendix - algorithms.tex}, lines 5--14}
{\src{BV}{1680--1688}; \src{S7}{329--395, 460--477};
 \src{STAT}{2324--2392}}

\finding{The BH family combines LD and DIF}
\man Step 2 states that p-values are adjusted for multiple testing.
The multiple-testing appendix explains BH without specifying its test family.

\dig The initial family is the list
\[
 \mathcal P=
 \bigl(p_{ii'}^{(i)}:i=1,\ldots,k,\ i'\ne i\bigr)
 \ \mathbin{\|}\
 \bigl(p_{ij}^{\mathrm{DIF}}:i=1,\ldots,k,\ j=1,\ldots,m\bigr),
\]
where \(\mathbin{\|}\) denotes concatenation. Its length is
\[
 T=k(k-1)+km.
\]
Both LD directions are included. The DIF contribution contains the single
statistic selected for each item--covariate pair, not both gamma and
chi-square p-values. Marginal diagnostics and subsequent score-effect tests
do not belong to this family. For AKQoL, \(k=10,m=2\), giving \(T=110\).

\effect Adjusting LD and DIF separately, adjusting one collapsed value per
unordered LD pair, or pooling all later tests produces a different procedure.
The manuscript does not establish which family generated its tables.
\trace{\path{article.tex}, lines 235 and 259;
 \path{Appendix - multiple testing.tex}, lines 1--6}
{\src{S7}{2965--2991}}

\finding{Cutoffs, unavailable tests, and boundary conventions}
For sorted values \(p_{(1)}\le\cdots\le p_{(T)}\), DIGRAM scans downward and
returns \(c_\alpha=r\alpha/T\) for the largest qualifying rank \(r\).
It returns \(\alpha/T\) if no rank qualifies. For ordinary valid p-values
this fallback still gives no rejections: failure at rank 1 implies
\(p_{(1)}>\alpha/T\). Thus the manuscript's rule of no rejection when no rank
qualifies is not a substantive disagreement.

The reported raw cutoff \(r\alpha/T\) can differ numerically from the largest
rejected observed p-value \(p_{(r)}\), while giving the same decisions on the
family used to compute it. Comparing BH-adjusted p-values with \(\alpha\)
is likewise equivalent for the same valid p-values and family.

DIGRAM initializes unavailable p-values with the sentinel \(2\), and its
collection loop retains these entries in \(T\). Dropping missing tests changes
the denominator. Its sorting routine also applies the floating-point
transformation \((1+p)-1\), which can affect exact boundary comparisons.
The source selects \(p\le c_\alpha\), whereas the manuscript sometimes says
that p-values must be below a threshold. Ties are allowed by DIGRAM; the
manuscript's strictly increasing ordering is already flagged by a TODO.
\trace{\path{article.tex}, line 259;
 \path{Appendix - multiple testing.tex}, lines 5--6}
{\src{MCA}{397--462}; \src{S7}{2380--2395, 2932--2940, 2965--2991}}

\finding{An estimate alone does not specify its p-value}
Algorithm 1 returns \(\widehat\gamma\), but later algorithms refer to it as
providing an asymptotic partial-gamma test and p-value. Neither that algorithm
nor the association-measures appendix specifies the variance estimate and
tail calculation needed to reproduce the test.

Initial source tests use \texttt{XYZ\_bias\_ANALYSE}. Its gamma result slot 6
is an asymptotic two-sided p-value, while slot 7 supplies the simulation-based
value when requested. Initial nominal DIF uses slots 3 and 4 for the
corresponding chi-square p-values. Later pruning uses a different result
structure, described in Section~\ref{sec:decisionp}.

This is a specification gap, not evidence that the manuscript's variance
formula is wrong: that formula is not provided. Handling an undefined gamma
is also incomplete. Algorithm 1 returns \(\mathrm{NA}\) when there are no
comparable pairs, without specifying the subsequent selection or multiplicity
decision. Section~\ref{sec:inference} records the source's numerical conventions.
\trace{\path{Appendix - algorithms.tex}, lines 59--70 and 284--295;
 \path{Appendix - association measures.tex}, line 15}
{\src{S7}{448--477, 808--815, 868--875}; \src{S3}{593--627}}

\section{Step 3(a): elimination of spurious LD evidence}
\label{sec:ld}

\finding{The weighted gamma agrees and remains fixed}
For \(h\in\{i,i'\}\), let
\[
 w_{ii'}^{(h)}
   =\sum_\ell\bigl(C_{\ell,ii'}^{(h)}+D_{\ell,ii'}^{(h)}\bigr).
\]
The manuscript's pair-level coefficient
\[
 \WPG_{ii'}=
 \frac{w_{ii'}^{(i)}\gamma_{ii'}^{(i)}
       +w_{ii'}^{(i')}\gamma_{ii'}^{(i')}}
      {w_{ii'}^{(i)}+w_{ii'}^{(i')}}
\]
agrees with the source's ratio of summed \texttt{IJXpmq} to summed
\texttt{IJXppq}, provided the denominator is positive and the participating
directional coefficients are defined, or zero-weight terms are omitted.
If one direction has no comparable pairs, the manuscript returns
\(\mathrm{NA}\), so a convention is needed for its zero-weight contribution;
the source stores gamma zero for that direction.
Both original directional
contributions remain in that ratio after one direction is suppressed.
Although Pascal reevaluates the ratio inside each search, the stored counts
are unchanged, so the weighted coefficient is effectively fixed.
The manuscript's initialization before the selection loop agrees with this.
\trace{\path{Appendix - algorithms.tex}, lines 88--123}
{\src{S7}{2455--2457, 2489--2490, 2524--2525, 2559--2560, 3293--3299}}

\finding{Original evidence and active evidence are different objects}
Let \(\mathcal H_0\) be the initial significant directional hypotheses, using
\(p\le c_{.05}\), and let \(\mathcal H_t\subseteq\mathcal H_0\) be those not
yet suppressed. For a pair, define
\begin{align*}
 N_+&=\sum_{h\in\{i,i'\}}
       \mathbf 1\{H_{ii'}^{(h)}\in\mathcal H_0,\ \gamma_{ii'}^{(h)}>0\},\\
 N_-&=\sum_{h\in\{i,i'\}}
       \mathbf 1\{H_{ii'}^{(h)}\in\mathcal H_0,\ \gamma_{ii'}^{(h)}\le0\},\\
 A_+(t)&=\sum_{h\in\{i,i'\}}
       \mathbf 1\{H_{ii'}^{(h)}\in\mathcal H_t,\ \gamma_{ii'}^{(h)}>0\},\\
 A_-(t)&=\sum_{h\in\{i,i'\}}
       \mathbf 1\{H_{ii'}^{(h)}\in\mathcal H_t,\ \gamma_{ii'}^{(h)}\le0\}.
\end{align*}
Pair subscripts are suppressed on these four counts for readability.
The source places a zero directional gamma in its nonpositive count;
its negative-stage names should be interpreted with this coding detail.

\dig \(N_+,N_-\) use the unchanged \texttt{part\_p} matrix.
\(A_+(t),A_-(t)\) use the temporary \texttt{temp\_p} matrix.

\man The appendix classifies pairs using only \(\mathcal H_t\), so an
originally two-direction pair is reclassified when one direction disappears.
\trace{\path{Appendix - algorithms.tex}, lines 128--179 and 233--236;
 \path{article.tex}, line 704}
{\src{S7}{2397--2426}}

\finding{The four source stages have different priorities}
\begin{center}
\begin{tabularx}{\linewidth}{@{}cL{0.27\linewidth}X@{}}
\toprule
Order & Source eligibility & Choice within the stage\\
\midrule
1 & \(N_+=2,\ A_+(t)>0\) & Largest strictly positive \(\WPG_{ii'}\).\\
2 & \(N_+=1,\ A_+(t)>0\) & Largest strictly positive \(\WPG_{ii'}\).\\
3 & \(N_-=2,\ A_-(t)>0\) & Smallest strictly negative \(\WPG_{ii'}\).\\
4 & \(N_-=1,\ A_-(t)>0\) & Smallest strictly negative \(\WPG_{ii'}\).\\
\bottomrule
\end{tabularx}
\end{center}

DIGRAM exhausts each stage before advancing to the next and never returns
to an earlier stage. It does not combine signs into a common two-direction
group and a common one-direction group. An originally two-positive pair
remains eligible in stage 1 after one direction is suppressed, because
the condition is \(A_+(t)>0\). The analogous statement holds in stage 3.

The manuscript instead gives any currently two-direction pair priority
over any currently one-direction pair and ranks within each group by
\(\lvert\WPG_{ii'}\rvert\). Both the sign priority and the use of original
counts differ from SCREEN J.
\trace{\path{article.tex}, lines 269--276 and 704;
 \path{Appendix - algorithms.tex}, lines 182--210}
{\src{S7}{2435--2584}}

\finding{Directional suppression itself agrees}
After selecting \(\{Y_a,Y_b\}\), the source sets every temporary p-value
in rows \(a\) and \(b\) to the sentinel \(2\). A row identifies the item
excluded from the rest-score. This is equivalent to the manuscript's update
\[
 \mathcal H_{t+1}
 =\mathcal H_t\setminus
 \{H_{ii'}^{(h)}:h\in\{a,b\}\}.
\]
It suppresses a directional test conditioned on \(R_a\) or \(R_b\).
It need not eliminate every unordered pair involving \(Y_a\) or \(Y_b\),
because the reverse direction may remain active.
No new statistical test is performed by suppression.
\trace{\path{article.tex}, line 276;
 \path{Appendix - algorithms.tex}, lines 217--236}
{\src{S7}{2472--2477, 2505--2510, 2541--2546, 2575--2580}}

\finding{Examples in which the retained LD evidence differs}
\textbf{Example 1: sign priority.}
Suppose only two unordered pairs have initial evidence.
Pair \(\{Y_1,Y_2\}\) has one significant positive direction, conditional on
\(R_1\), with \(\WPG_{12}=0.40\).
Pair \(\{Y_1,Y_3\}\) has two significant negative directions, with
\(\WPG_{13}=-0.80\). All other directions are nonsignificant.

SCREEN J selects \(\{Y_1,Y_2\}\) in its one-positive stage.
Suppressing rows 1 and 2 leaves the \(R_3\) direction of
\(\{Y_1,Y_3\}\), which remains eligible in the original two-negative stage.
Both pairs are selected as LD evidence.
The manuscript selects \(\{Y_1,Y_3\}\) first, suppresses row 1, and loses
the only significant direction of \(\{Y_1,Y_2\}\).
This example concerns evidence selection before the separate negative-LD
model-reporting policy in Section~\ref{sec:negative}.

\textbf{Example 2: original counts versus reclassification.}
Suppose pairs \(\{Y_1,Y_2\}\), \(\{Y_1,Y_3\}\), and
\(\{Y_3,Y_4\}\) each have two significant positive directions, with weighted
gammas \(0.90,0.80,0.70\), respectively. No other pair has evidence.
Both procedures first select \(\{Y_1,Y_2\}\).

SCREEN J still treats \(\{Y_1,Y_3\}\) as an originally two-positive pair
and selects it next. The remaining \(R_4\) direction then permits selection
of \(\{Y_3,Y_4\}\). The manuscript reclassifies \(\{Y_1,Y_3\}\) as a
one-direction pair, selects \(\{Y_3,Y_4\}\) next, and suppresses the last
active direction of \(\{Y_1,Y_3\}\).
The selected sets are
\[
 \mathcal G_{\mathrm{DIGRAM}}=\{12,13,34\},\qquad
 \mathcal G_{\mathrm{manuscript}}=\{12,34\},
\]
where \(ii'\) abbreviates \(\{Y_i,Y_{i'}\}\).
These hypothetical summaries isolate control-flow differences; they are
not asserted to be the output of newly fitted datasets.

\finding{Mixed signs, ties, and termination}
The manuscript's four active lists exclude a pair with two significant
directions of opposite sign. Such a pair is neither two-positive,
two-negative, nor exactly one significant direction. If only such pairs
remain, its loop has active hypotheses but no eligible maximizer.

The source counts signs separately. A mixed-sign pair has \(N_+=N_-=1\);
its active sign and the sign of its fixed pooled gamma determine eligibility
in the corresponding stage. Positive stages require \(\WPG>0\); negative
stages require \(\WPG<0\). Maximizing absolute gamma without these conditions
is different. These cases are not claimed to occur in the manuscript's
data examples.

Opposite directional signs are possible even with positive marginal
item associations. As a concrete counting illustration, consider three
binary items with the following frequencies:
\begin{center}
\begin{tabular}{@{}cccccccc@{}}
\toprule
\(000\)&\(001\)&\(010\)&\(011\)&\(100\)&\(101\)&\(110\)&\(111\)\\
\midrule
20&2&1&3&3&1&2&20\\
\bottomrule
\end{tabular}
\end{center}
The pattern digits are \(Y_1Y_2Y_3\). Direct pair counting gives
\[
 \gamma_{12}^{(1)}=\frac{4-1}{4+1}=0.6,\qquad
 \gamma_{12}^{(2)}=\frac{4-9}{4+9}=-\frac5{13},\qquad
 \WPG_{12}=-\frac19.
\]
All three marginal item-pair gammas are positive.
Scaling the frequencies preserves these coefficients and can make both
nonzero directional associations asymptotically significant.
The missing mixed-sign case is therefore not merely an impossible input.

The source scans \(i<i'\) in item order and updates the best pair only on a
strict improvement, retaining the first exact tie. The manuscript's
\(\arg\max\) does not specify tie-breaking.
DIGRAM exits after exhausting its four eligible stages, which need not mean
all temporary significant entries have been suppressed.
Its coefficient assignments are guarded by a positive denominator; this is
not a fully specified recovery rule for an undefined weighted gamma.
\trace{\path{Appendix - algorithms.tex}, lines 126--210}
{\src{S7}{2402--2425, 2435--2439, 2448--2458, 2487--2491,
 2522--2526, 2557--2584}; \src{STAT}{1004--1030}}

\section{Steps 3(b) and 3(c): DIF purification}
\label{sec:dif}
\finding{The passes share a changing DIF matrix}
Let \(\mathcal E_0\subseteq\Iset\times\{X_1,\ldots,X_m\}\) be the initial
DIF relations. The manuscript runs two independent reductions of
\(\mathcal E_0\), then retains their intersection. It explicitly states
that their order does not matter.

DIGRAM uses the following sequence:
\begin{enumerate}
\item For each \(X_j\), form its current item set \(\DIF(X_j)\),
apply the all-items exception below if needed, and run the per-covariate
item-pruning procedure when applicable: manuscript Step 3(c).
\item Write each reduced set directly into the shared DIF matrix.
\item For each \(Y_i\), form \(\SRC(Y_i)\) from that already reduced
matrix and run the per-item source-pruning procedure when applicable:
manuscript Step 3(b).
\item Write these reductions into the same matrix.
The outer pair of passes is not iterated to convergence.
\end{enumerate}
Let \(\mathsf P_b^{D},\mathsf P_c^{D}\) denote the source's two reductions,
including its all-items handling in \(\mathsf P_c^{D}\), and let
\(\mathsf P_b^{M},\mathsf P_c^{M}\) denote the manuscript's reductions.
The resulting relations are
\[
 \mathcal E_{\mathrm{DIGRAM}}
 =\mathsf P_b^{D}\bigl(\mathsf P_c^{D}(\mathcal E_0)\bigr),\qquad
 \mathcal E_{\mathrm{manuscript}}
 =\mathsf P_b^{M}(\mathcal E_0)\cap\mathsf P_c^{M}(\mathcal E_0).
\]
Even if the individual reduction rules were matched, composing the passes
need not equal intersecting their independent results.
Deleting a relation changes the conditioning sets available to the
subsequent pass. The statistic and all-items exceptions described below
also make the individual source and manuscript reductions different.
\trace{\path{article.tex}, lines 279--285 and 709}
{\src{S7}{3036--3105, 3116--3155}}

\finding{The hypotheses agree for identical candidate sets}
For a fixed per-covariate set \(\mathcal D_j\), the first source pass tests
\[
 Y_i\ind X_j\mid S,\ \mathcal D_j\setminus\{Y_i\},
 \qquad Y_i\in\mathcal D_j.
\]
For a fixed per-item source set \(\mathcal S_i\), the second tests
\[
 Y_i\ind X_j\mid S,\ \mathcal S_i\setminus\{X_j\},
 \qquad X_j\in\mathcal S_i.
\]
These are the forms described for manuscript Steps 3(c) and 3(b).
The differences concern which sets enter the tests and the decision p-value.

Within each invoked purification routine, all tests for the current set
are performed before at most one candidate is removed.
The source removes the candidate with the largest decision p-value if it
exceeds \(0.05\), rebuilds the table, and repeats.
The manuscript's one-at-a-time removal after a complete pass agrees.
For equal largest decision p-values, the source chooses the first candidate
in source order because its running maximum changes only on a strict
increase. The manuscript does not specify a tie-breaking rule.
\trace{\path{article.tex}, lines 279--283;
 \path{Appendix - algorithms.tex}, lines 267--321}
{\src{S13}{686--773, 1108--1195}}

\finding{The decision p-value combines two tests}
\label{sec:decisionp}
Write \(p_\chi\) for the configured contingency-table p-value and
\(p_{\gamma,2}\) for the two-sided gamma p-value.
In asymptotic mode, both source purification routines use
\[
 p_{\mathrm{decision}}=\min\{p_\chi,p_{\gamma,2}\}.
\]
The follow-up result slot 6 is a one-tail gamma probability, so the code
forms \(p_{\gamma,2}=2\,\texttt{Results}[6]\).
This differs from the already two-sided slot 6 of the initial
\texttt{XYZ\_bias\_ANALYSE} results.
In simulation mode, the follow-up rule takes the minimum of slots 4 and 9.

The manuscript instead describes the partial-gamma p-value as the decision
value. For example, \(p_{\gamma,2}=0.20\) and \(p_\chi=0.01\) give a source
decision value of \(0.01\). Gamma-only elimination could discard the
candidate. The minimum-p rule can also change which candidate is removed
first and thus all later conditioning sets.

These removal thresholds are \(0.05\), not the initial \(c_{.05}\).
BH values printed for follow-up diagnostics do not replace the explicit
\(p_{\max}>0.05\) pruning condition. Stage-specific type handling and
suppressed gamma calculations are discussed in Section~\ref{sec:inference}.
\trace{\path{article.tex}, line 280;
 \path{Appendix - algorithms.tex}, lines 284--295}
{\src{S13}{732--773, 1154--1195};
 \src{XYZ}{805--814}; \src{S3}{614--627}}

\finding{The all-items-DIF exception changes the candidate set}
For a covariate with \(\lvert\DIF(X_j)\rvert=k\) initially, DIGRAM sets
\[
 \DIF(X_j)\leftarrow
 \{Y_i:p_{ij}^{\mathrm{DIF}}\le c_{.01}\}.
\]
This uses the \(1\%\) cutoff from the same initial global BH family,
not an unadjusted \(0.01\) rule or a newly fitted test.

If all \(k\) items still survive, the source warns and bypasses ordinary
per-covariate item purification for that covariate.
A set with two to \(k-1\) members is purified. An empty set or incoming
singleton does not trigger the routine.
The resulting matrix still feeds the subsequent per-item source pass.
No corresponding exception appears in the manuscript's Step 3(c) description.
\trace{\path{article.tex}, lines 282--285}
{\src{S7}{2391--2395, 3053--3103}}

\finding{Singletons do not universally receive new purification tests}
SCREEN J invokes either purification routine only when its incoming list
has more than one candidate, subject to the all-items exception.
An incoming singleton is carried forward without a new call.
Thus not every DIF pair receives two new purification p-values.

This is an entry condition, not a rule protecting the final member of an
invoked DIF routine. Once purification has started, its internal loop can
retest a reduced singleton and remove it.
The score-effect routine has a different stopping rule
(Section~\ref{sec:step4}).
\trace{\path{article.tex}, lines 285 and 615--627;
 \path{Appendix - data example.tex}, lines 55--62}
{\src{S7}{3081--3101, 3133--3153}; \src{S13}{491--773, 944--1195}}

\finding{The hypothetical DIF example follows a different test history}
The manuscript starts with
\[
 \SRC(A)=\{X,Z\},\quad \SRC(C)=\{X\},\quad \SRC(E)=\{Z\},
\]
or \(\DIF(X)=\{A,C\}\), \(\DIF(Z)=\{A,E\}\).
Its illustrated per-covariate pass removes C--X and A--Z, leaving A--X
and E--Z.

Conditional on those illustrated decisions, SCREEN J's next per-item pass
sees singleton source lists. It would not execute the manuscript's
earlier illustrated tests
\[
 A\ind X\mid S,Z,\qquad A\ind Z\mid S,X.
\]
The final edge set can happen to agree while the test history differs.
This comparison holds the illustrated decisions fixed to isolate the
pass-order issue; DIGRAM's minimum-p statistic could itself change them.
\trace{\path{article.tex}, lines 289--401}
{\src{S7}{3088--3093, 3124--3153}}

\section{Step 4: associations between score and covariates}
\label{sec:step4}
\finding{Marginal preselection precedes conditional elimination}
\man The algorithm begins with the candidate covariate set \(\mathcal C\)
and tests \(S\ind X_j\mid\mathcal C\setminus\{X_j\}\) for every member.
No preliminary marginal selection rule is defined.

\dig The source clears its selected score-effect vector, scans all
exogenous variables, and tests each marginal association \(S\ind X_j\).
It initializes
\[
 \mathcal C_0=\{X_j:
 \min(p_{\chi,j}^{\mathrm{marg}},p_{\gamma,2,j}^{\mathrm{marg}})\le0.05\}.
\]
Only members of \(\mathcal C_0\) enter conditional backward elimination.

\effect A covariate whose association appears only after adjustment can
enter and survive the manuscript's procedure but fail DIGRAM's marginal
filter. Changing the initial conditioning set can also alter tests for
other covariates. An older \(0.01\) option is commented out in the source;
the examined implementation uses \(0.05\).
\trace{\path{article.tex}, lines 407--410;
 \path{Appendix - algorithms.tex}, lines 340--355}
{\src{S13}{1335--1346, 1382--1395, 1527--1529, 1568--1588}}

\finding{The statistic and stopping rule differ}
With more than one selected effect, DIGRAM tests
\[
 S\ind X_j\mid\mathcal C_t\setminus\{X_j\},\qquad X_j\in\mathcal C_t,
\]
and removes the candidate with the largest
\(\min(p_\chi,p_{\gamma,2})\) when that value exceeds \(0.05\).
This is separate from the initial LD/DIF BH family.

The manuscript uses gamma alone and loops while \(\mathcal C\ne\varnothing\).
DIGRAM stops conditional elimination when at most one candidate remains.
The singleton has already passed marginal selection.
A final conditional-loop singleton retest, such as the last row of the
AKQoL Step 4 table, is not a separate source pruning iteration.
When an elimination leaves one candidate, its stored score--covariate
gamma remains the value from the last conditional table; the source does
not recompute a final marginal gamma. Equal largest decision p-values
are resolved by choosing the first candidate in source order.

The hypotheses and removal of at most one candidate per completed pass
otherwise agree for identical input sets.
Marginal and conditional tables have their own record-selection rules,
so identical observations cannot be assumed when covariates have missing values.
\trace{\path{Appendix - algorithms.tex}, lines 347--366;
 \path{article.tex}, lines 631--645}
{\src{S13}{1660--1664, 1858--1863, 1883--1929}}

\section{Step 5: graph construction and treatment of negative LD}
\label{sec:step5}
\finding{SCREEN J reports evidence without replacing the current model}
The manuscript assembles a graphical log-linear Rasch model (GLLRM),
an item response theory (IRT) graph, and a moralized marginal graph.
SCREEN J itself does not perform that construction. Its command handler
reports positive LD, negative LD, DIF, and score associations, then restores
the previously saved LD and DIF matrices.

The related non-J branch initializes graph structures and defines a model.
This establishes the scope difference; it does not establish that the
manuscript's graph-building code reproduces every detail of that separate
branch. The ordinal-model formulas in the manuscript's appendix describe
a statistical model, while SCREEN J is a screening operation and does not
fit its interaction parameters.
\trace{\path{article.tex}, lines 412--417 and 649--665;
 \path{Appendix - general item response model.tex}, lines 21--32}
{\src{CMD}{1793--1804, 1809--1922}; \src{S7}{3286--3301}}

\finding{Selected negative LD is excluded from model reporting}
\label{sec:negative}
For a selected pair, the source command handler uses the sign of the
sum of the initial cached directional gammas,
\[
 \gamma_{ii'}^{(i)}+\gamma_{ii'}^{(i')}
\]
to distinguish positive from nonpositive LD. If that sum is nonpositive,
the relation is put in the negative-LD report with the statement that it
is not included in the model. This classification uses the unweighted sum;
Step 3(a) ranks pairs using the weighted gamma.

The manuscript retains selected negative associations as model edges.
Q3--Q10 has initial weighted gamma \(-0.694\), appears in its tentative
and final GLLRM, and subsequently has two negative Step 6 estimates,
\(-0.611,-0.755\). The latter are final graph checks and must not be
substituted for the initial directional values in DIGRAM's reporting rule.
The manuscript gives only one initial directional estimate for this pair.
Determining its exact source classification requires both initial values;
a negative weighted gamma alone does not establish that their unweighted
sum is nonpositive.

The treatment of evidence of negative LD and a retained model interaction
is therefore a substantive policy difference. Because SCREEN J ultimately
restores the current model, this does not mean it directly edits a final
fitted GLLRM.
\trace{\path{article.tex}, lines 600--611, 650--665, and 690--697}
{\src{CMD}{1818--1853, 1918--1922}}

\section{Step 6: comparison with the separate graph checks}
\label{sec:step6}
\finding{There is no final C5/C7/C9 stage in SCREEN J}
After DIF and score screening, SCREEN J reports warnings when an LD pair's
items share a DIF source, caches partial statistics, and finishes.
It does not apply the manuscript's rule of retaining an LD edge when
at least one final direction is significant.
The AKQoL Step 6 removal of Alder--Q4 is not a SCREEN J operation.

For the remaining comparisons in this section, the reference is explicitly
DIGRAM's separate CHECK I machinery. CHECK I invokes LD checks for present
edges and local-independence checks for absent edges, along with checks of
retained DIF and no-DIF assumptions. It records and reports statistics.
Its active loop does not delete model edges under the manuscript's rule.
\trace{\path{article.tex}, lines 420--428 and 669--711}
{\src{S7}{3233--3301}; \src{CMD}{3382--3393};
 \src{EXE}{4577--4646}}

\finding{LD-check DIF sources come from the other side of the component split}
Remove \(\{Y_a,Y_b\}\), and let \(C_a,C_b\) be the remaining LD components
containing its endpoints. For distinct components, the manuscript's
explanation corresponds to
\begin{align}
 Y_a&\ind Y_b\mid S_{C_a^c},\ \SRC(C_a),\label{eq:ld-man-a}\\
 Y_a&\ind Y_b\mid S_{C_b^c},\ \SRC(C_b).\label{eq:ld-man-b}
\end{align}
The manuscript's displayed \(\DIF(C_a)\) is interpreted here according
to its prose as sources of DIF for the component's items.

The source's \texttt{remove\_component} first removes the component's items
from the score indicator vector, then collects sources of DIF for items
\emph{still included in that score}. Its tests are
\begin{align}
 Y_a&\ind Y_b\mid S_{C_a^c},\ \SRC(C_a^c),\label{eq:ld-dig-a}\\
 Y_a&\ind Y_b\mid S_{C_b^c},\ \SRC(C_b^c).\label{eq:ld-dig-b}
\end{align}
The subscore definitions agree; the additional DIF covariate sets differ.
Renaming the overloaded \(\DIF(C)\) alone would not resolve this.
\trace{\path{article.tex}, lines 428 and 488--498}
{\src{S13}{2350--2363, 2512--2527, 2850--2882}}

In the A--B example,
\[
 C_A=\{A,D,C\},\quad C_B=\{B\},\quad
 \SRC(C_A)=\{X\},\quad \SRC(C_B)=\{X,Z\}.
\]
\begin{center}
\begin{tabularx}{\linewidth}{@{}L{0.24\linewidth}XX@{}}
\toprule
Excluded component & Manuscript conditioning & DIGRAM LDCheck conditioning\\
\midrule
\(\{A,D,C\}\) & \(S_{\{B\}},X\) & \(S_{\{B\}},X,Z\)\\
\(\{B\}\) & \(S_{\{A,D,C\}},X,Z\) & \(S_{\{A,D,C\}},X\)\\
\bottomrule
\end{tabularx}
\end{center}
This table compares the stated sets; the first row also has a degeneracy
discussed in Section~\ref{sec:exampleissues}.

\finding{The source DIF check conditions on all other biased items}
The manuscript's C5 hypothesis is
\begin{equation}
 Y_i\ind X_j\mid
 S,\ \DIF(X_j)\setminus\{Y_i\},\
 \SRC(Y_i)\setminus\{X_j\}.
 \label{eq:dif-man}
\end{equation}
DIGRAM's separate \texttt{DIFCheck} marks every other item having any retained
DIF source. It forms a common table with those items, the tested item's
sources, and the score, then conditions on all variables except the tested
item--source pair. In the notation of Section~\ref{sec:notation}, this is
\begin{equation}
 Y_i\ind X_j\mid
 S,\ \mathcal B_{\mathrm{all}}\setminus\{Y_i\},\
 \SRC(Y_i)\setminus\{X_j\}.
 \label{eq:dif-dig}
\end{equation}
The source part agrees, but the conditioned item set can be larger in DIGRAM:
an item having DIF from another covariate is still included.

For the tentative AKQoL structure, the source would check
\[
 Q4\ind\texttt{Alder}\mid S,Q5,\qquad
 Q5\ind\texttt{woman}\mid S,Q4.
\]
The manuscript's C5 formula conditions only on \(S\) for these relations.
Similarly, the hypothetical check of B against Z includes C in DIGRAM
because C has DIF from X, giving \(B\ind Z\mid S,C,X\).
The manuscript's table instead gives \(B\ind Z\mid S,X\).
These are different hypotheses, not replacement numerical results.
\trace{\path{article.tex}, lines 421--426, 441, 616--624, and 675--678}
{\src{S13}{3736--3803, 3940--3945, 4015--4021, 4064--4072}}

\finding{General graph checks have additional guards and data restrictions}
After removing the tested LD edge, \texttt{LDCheck} skips its component
tests when \(C_a=C_b\), as occurs for an edge on an LD cycle.
The manuscript's general instruction to test every retained pair has
no such exception.

This guard can be irrelevant to a graph obtained solely by SCREEN J's
LD-selection rule: each new selection requires at least one unsuppressed
endpoint row and therefore introduces a previously unused endpoint.
Such selection cannot close an LD cycle.
The guard remains relevant for a general supplied graph and is not asserted
to cause a discrepancy in the data examples.

Both source graph-check routines use \texttt{GET\_GOOD\_ITEMS} and therefore
require a valid, non-extreme total score in the configured range.
The LD test additionally keeps only
\[
 0<S_{C_a^c}<|C_a^c|\,\texttt{item\_max}
\]
in the first direction, with the analogous rule in the second.
The upper bound is written as implemented; it is not silently replaced by
a sum of heterogeneous item maxima. These restrictions are absent from the
manuscript's general Step 6 description.
\trace{\path{article.tex}, line 428}
{\src{S13}{2485--2487, 2601--2612, 2809--2847, 3862--3872};
 \src{S2}{3388--3403}}

\finding{The SPADI pain checks are not redundant under the source rule}
The pain example retains P3--P4 LD and P1--\texttt{over60} DIF.
Removing P3--P4 leaves components \(\{P3\}\) and \(\{P4\}\).
Both complementary subscores contain P1. The source LD checks therefore
additionally condition on \texttt{over60}:
\[
 P3\ind P4\mid S-P3,\ \texttt{over60},\qquad
 P3\ind P4\mid S-P4,\ \texttt{over60}.
\]
The initial Step 2 tests condition only on the respective rest-score.
These are different hypotheses, so DIGRAM's rule does not justify the
manuscript's statement that the final checks are redundant for both subscales.

For disability D4--D5, the example has no retained DIF. This particular
conditioning-set objection does not apply there.
Even matching symbolic hypotheses would not guarantee identical numeric
tests without matching record filters and inference conventions.
\trace{\path{Appendix - data example.tex}, lines 19, 47, 62, and 308}
{\src{S13}{2350--2363, 2604--2612, 2850--2882}}

\section{Internal example and reproducibility issues}
\label{sec:exampleissues}
The following findings concern consistency within the manuscript.
They should be distinguished from direct disagreement with SCREEN J.

\finding{An adjusted p-value above 0.05 is carried forward}
The AKQoL initial DIF table gives Q4--Alder raw \(p=0.0031\) and
\(p_{\mathrm{adj}}=0.062\), but the text counts it among the two flagged
DIF relations and carries it forward.
If the displayed adjusted p-value is the decision value under the stated
\(5\%\) criterion, that decision does not follow.

The source comparison does not identify the intended correction.
The adjustment family, candidate-generation threshold, or table needs
reconciliation. The displayed values alone do not identify the adjustment
method that was actually used.
\trace{\path{article.tex}, lines 259, 568--580, and 615--627}
{\src{S7}{2385--2389, 2965--2991}}

\finding{The same stated DIF hypothesis has different reported p-values}
Each of Q4--Alder and Q5--\texttt{woman} is the only DIF relation for its
item and source in the tentative AKQoL model.
The manuscript's C5 formula therefore reduces to
\[
 Q4\ind\texttt{Alder}\mid S,\qquad
 Q5\ind\texttt{woman}\mid S,
\]
the initial DIF hypotheses.
The raw p-values nevertheless change from \(0.0031,0.0008\) to
\(0.051,0.032\), while the gammas remain \(0.383,0.727\).

With the same observations, test, variance, tail convention, and adjustment
status, a repeated hypothesis should reproduce the same p-value.
A change in any of those ingredients could explain the different values,
but it must be stated. Equal gamma estimates alone do not prove that the
inferential calculations were identical.
DIGRAM's larger item-conditioning set in \eqref{eq:dif-dig} defines another
possible test; it does not validate these particular numbers.
\trace{\path{article.tex}, lines 421--426, 623--627, and 675--681}
{\src{S13}{3736--3803, 4064--4072}}

\finding{The hypothetical Step 6 example has two distinct problems}
First, the LD table implies retaining A--B and C--D and removing A--D under
the manuscript's at-least-one-significant-direction rule.
The prose and final graph instead retain A--D and remove C--D.
This contradiction is already identified by the margin TODO.

Second, the example defines \(C_A^c=\{B\}\), hence \(S_{C_A^c}=B\).
Its first purported LD test conditions on B itself:
\[
 A\ind B\mid B,X.
\]
This conditional independence is automatic. B is constant within every
conditioning stratum, so there are no comparable pairs for a nontrivial
partial-gamma test. The significant p-value \(0.016\) cannot represent
that displayed partial-gamma test. Adding Z, as the source's LD-check rule
requires, does not remove the degeneracy.

For binary items with \(\texttt{item\_max}=1\), DIGRAM's additional
exclusion of extreme complementary subscores would also leave no
observations for this direction.
For ordinal B, conditioning on B still removes its within-stratum variation.
Thus the example needs a check of its test setup as well as its final edges.
\trace{\path{article.tex}, lines 488--522}
{\src{S13}{2604--2612}; the degeneracy also follows directly from the
displayed manuscript hypothesis}

\finding{The subscore count needs a precise definition}
The manuscript states that excluding the item under consideration and the
zero subscore leaves \(2^{k-1}\) subscores.
There are \(2^{k-1}-1\) nonempty subsets of the remaining \(k-1\) items.
Distinct subsets can also give the same realized numerical scores.
The count should therefore specify whether it concerns subset definitions
or distinct observed score variables.
This issue is already marked by a TODO and is separate from the source's
non-exhaustive screening scope.
\trace{\path{article.tex}, lines 701--702}
{\src{S7}{2943--3008}; Section~\ref{sec:step1} of this note}

\section{Inference and implementation details for numerical reproduction}
\label{sec:inference}
The manuscript explicitly chooses asymptotic inference.
The existence of a simulation option in DIGRAM is not itself a discrepancy.
However, matching that choice does not match all numerical conventions.
The details below extend the step-by-step comparison where a generic
instruction to compute gamma or a p-value is insufficient.

\finding{Result slots and tails differ between initial and follow-up tests}
\begin{center}
\begin{tabularx}{\linewidth}{@{}L{0.29\linewidth}XX@{}}
\toprule
Decision stage & Asymptotic source value & Simulation source value\\
\midrule
Initial gamma & \(\texttt{Bias}[6]\), already two-sided &
\(\texttt{Bias}[7]\)\\
Initial nominal DIF & \(\texttt{Bias}[3]\) & \(\texttt{Bias}[4]\)\\
DIF purification and score screening &
\(\min(\texttt{BT}[3],\,2\texttt{BT}[6])\) &
\(\min(\texttt{BT}[4],\,\texttt{BT}[9])\)\\
\bottomrule
\end{tabularx}
\end{center}
Here \texttt{Bias} denotes the initial \texttt{XYZ\_bias\_ANALYSE}
result structure and \texttt{BT} the follow-up big-table result structure.
The slot numbers identify implementation outputs, not model parameters.
The competing \(p_\chi\) comes from the configured contingency-table test.
Within conditioning strata, its statistics and degrees of freedom are
summed before evaluating the pooled asymptotic chi-square tail.
\trace{\path{article.tex}, line 235;
 \path{Appendix - algorithms.tex}, lines 284--295 and 353--355}
{\src{S3}{483--503, 614--627}; \src{XYZ}{805--814, 1693--1717};
 \src{S13}{732--750, 1154--1172, 1883--1901};
 \src{STAT}{2324--2392}}

\finding{The source variance and normal-tail conventions are explicit}
In stratum \(\ell\), let \(n_{uv,\ell}\) be the count in cell \((u,v)\)
and \(n_\ell=\sum_{uv}n_{uv,\ell}\).
Let \(a_{uv,\ell}\) and \(d_{uv,\ell}\) count observations concordant and
discordant, respectively, with a response in that cell. Then the source uses
\[
 \mathrm{PMQ}_\ell=\sum_{uv}n_{uv,\ell}(a_{uv,\ell}-d_{uv,\ell}),\quad
 \mathrm{PPQ}_\ell=\sum_{uv}n_{uv,\ell}(a_{uv,\ell}+d_{uv,\ell}),
\]
which equal \(2(C_\ell-D_\ell)\) and \(2(C_\ell+D_\ell)\).
For a nonempty stratum with comparable pairs, its variance numerator is
\[
 v_\ell=4\left[
  \sum_{uv}n_{uv,\ell}(a_{uv,\ell}-d_{uv,\ell})^2
  -\frac{\mathrm{PMQ}_\ell^2}{n_\ell}
 \right].
\]
The partial routines pool these terms and use
\[
 \widehat V(\widehat\gamma)
 =\frac{\sum_\ell v_\ell}{(\sum_\ell\mathrm{PPQ}_\ell)^2}.
\]
Ordinarily, the test uses
\(U=|\widehat\gamma|/\sqrt{\widehat V(\widehat\gamma)}\).
If the pooled denominator is positive but the variance is nonpositive,
the code instead assigns \(U=4\).
If \(U>4\), it assigns the one-tail probability \(10^{-6}\), rather than
evaluating the normal tail beyond 4. Otherwise it calls its normal-tail
routine. The two-sided decision probability is formed as described above.

These are preserved implementation conventions, not recommendations for
a new inference method. A conventional unmodified normal-tail calculation
can disagree numerically even when the gamma coefficient matches.
The manuscript does not specify the variance or these boundary conventions.
\trace{\path{Appendix - algorithms.tex}, lines 63--70;
 \path{Appendix - association measures.tex}, lines 7--15}
{\src{STAT}{1004--1030, 1296--1334};
 \src{S3}{486--522}; \src{XYZ}{786--814}}

\finding{Undefined and unavailable tests use sentinels, not NA}
When pooled comparable-pair counts are zero, the initial partial-gamma
routine has gamma initialized to zero and assigns a one-tail probability
of 1. Its stored two-sided value is therefore 2, without truncation to 1.
This is a sentinel-like non-rejection value, not a valid ordinary p-value.
The manuscript instead returns \(\mathrm{NA}\) and does not define its
subsequent treatment.

In follow-up \texttt{XYZ\_TEST}, zero effective sample size or zero total
degrees of freedom invokes \texttt{NoResults}, setting gamma to 0 and
result slots 3, 4, 6, and 7 to 2.
The asymptotic pruning comparison consequently gives
\(\min(2,2\cdot2)=2\), hence nonsignificance.
An empty or oversized multidimensional table can instead cause an earlier
routine exit, so not every no-data event is an ordinary elimination step.

An exact-mode boundary requires further caller-state inspection:
\texttt{NoResults} does not assign slot 9, although exact purification
normally reads it. This note records that omission without claiming a
particular resulting exact p-value or demonstrating a runtime error.
The manuscript's asymptotic scope does not require resolving that
simulation-only boundary.
\trace{\path{Appendix - algorithms.tex}, lines 30--37 and 59--61}
{\src{S3}{447--448, 508--522, 619--627};
 \src{XYZ}{960--975, 1697--1712};
 \src{S13}{567--568, 680--681, 732--750}}

\finding{Variable types are handled differently across source stages}
The source does not follow a single universal nominal/ordinal rule:
\begin{enumerate}
\item Initial DIF selects gamma for binary or ordinal covariates and
the contingency-table statistic otherwise.
\item Per-covariate DIF-item purification retains the exogenous variable's
original type and assigns items type 3. For a multilevel nominal covariate,
the gamma test is suppressed, so in asymptotic mode the minimum-p decision
effectively uses the contingency-table test. In simulation mode, the
unassigned-slot-9 qualification above also matters when gamma is suppressed.
\item Per-item multiple-source purification assigns every source type 3
in the tested table, although it separately remembers original types for
reporting.
\item Both marginal and conditional score screening assign their exogenous
variables type 3.
\end{enumerate}
Thus the latter stages can calculate an order-dependent gamma even for a
covariate originally marked nominal. This describes the preserved source;
it is not an endorsement of assigning an order to nominal categories.
The manuscript's gamma-only formulation is different, but replacing it
with an unconditional nominal-always-chi-square description would still
not reproduce these source paths.
\trace{\path{article.tex}, lines 235 and 279--285;
 \path{Appendix - algorithms.tex}, line 353}
{\src{BV}{1680--1688}; \src{S7}{460--477};
 \src{S13}{515--540, 953--964, 1388--1394, 1694--1704};
 \src{XYZ}{751--755, 1716--1717}}

\finding{Tables can use different observations after a candidate is removed}
The per-covariate DIF pass requires valid item responses, a non-extreme
allowed total score, and the target covariate's value.
The per-item multiple-source pass requires values of all sources in its
current candidate set, including the source being tested. It rebuilds the
table after a removal.
Dropping a source with missing values can therefore admit observations
that were previously excluded.

Score screening likewise uses one source's valid values for a marginal
table but all retained sources' values for a conditional table.
It includes extreme total scores.
There is also a literal capacity-boundary difference: its marginal pass
caps scores \(S\ge\texttt{maxdim}\) at \(\texttt{maxdim}-1\), whereas the
conditional pass excludes \(S\ge\texttt{maxdim}\).
These details can be irrelevant in a small complete dataset, but they must
not be silently replaced by a universal fixed-sample assumption.
\trace{\path{article.tex}, line 562;
 \path{Appendix - algorithms.tex}, lines 13--37 and 343--355}
{\src{S13}{620--642, 1050--1069, 1458--1479, 1765--1784}}

\finding{Simulation inference adds further matching requirements}
The manuscript uses asymptotic tests, so this is a reproducibility extension
for any later comparison using DIGRAM's exact or repeated-simulation modes.
The source counts simulated contingency-table statistics at least as large
as the observed statistic, and absolute simulated gammas at least as large
as the observed absolute gamma. It divides exceedance counts by simulations
performed, without an add-one correction.
Follow-up result slot 7 is directional; slot 9 is the two-sided
absolute-gamma result read by purification.

Matching simulation output requires the simulation count, generator and
random state, table generation, and sequential-stopping settings.
A change in those settings is not a demonstrated disagreement in the
manuscript's asymptotic examples.
\trace{\path{article.tex}, line 235}
{\src{S3}{121--145}; \src{XYZ}{839--865, 1721--1789};
 \src{S13}{736--745, 1158--1167}}

\section{Source catalog}
\label{sec:sources}
All original-source paths below are relative to
\path{../digram/source/digram_source_20260817/}.
Line numbers refer to the baseline in Section~\ref{sec:scope}.
Printed paths and line numbers keep this document usable on Overleaf
without requiring the sister repository as a compilation dependency.

\begin{longtable}{@{}L{0.09\linewidth}L{0.32\linewidth}L{0.52\linewidth}@{}}
\toprule
Key & Original file & Relevant role\\
\midrule
\endfirsthead
\toprule
Key & Original file & Relevant role\\
\midrule
\endhead
\bottomrule
\endfoot
\hypertarget{source-CMD}{CMD} & \path{scd/DIGRAM1f.pas} &
SCREEN J dispatch, saved model state, negative-LD reporting,
graph construction, and CHECK command dispatch.\\
\hypertarget{source-EXE}{EXE} & \path{skunits/DGRexe.pas} &
Screening wrapper and separate CHECK I execution loop.\\
\hypertarget{source-S7}{S7} & \path{skunits/SKbias7.pas} &
Initial LD/DIF tests, combined BH family, LD selection,
DIF-pass orchestration, and screening summaries.\\
\hypertarget{source-S13}{S13} & \path{skunits/SKbias13.pas} &
DIF purification, score-effect selection, and separate LDCheck
and DIFCheck routines.\\
\hypertarget{source-S2}{S2} & \path{skunits/SKbias2.pas} &
Item extraction, reversal, and valid-score filtering.\\
\hypertarget{source-S3}{S3} & \path{skunits/SKbias3.pas} &
Initial conditional test calculations and result slots.\\
\hypertarget{source-MCA}{MCA} & \path{skunits/SKmca.pas} &
BH sorting and raw-cutoff computation.\\
\hypertarget{source-BV}{BV} & \path{skunits/BIASvars.pas} &
Covariate-type interpretation, including the binary-variable rule.\\
\hypertarget{source-STAT}{STAT} & \path{skunits/SkStat.pas} &
Concordant/discordant counts, marginal gamma inference,
and configured contingency-table statistics.\\
\hypertarget{source-XYZ}{XYZ} & \path{skunits/SKxyz1.PAS} &
Follow-up conditional tests and simulation conventions.\\
\end{longtable}

Manuscript references concern \path{article.tex},
\path{Appendix - algorithms.tex},
\path{Appendix - association measures.tex},
\path{Appendix - multiple testing.tex},
\path{Appendix - general item response model.tex}, and
\path{Appendix - data example.tex}.
Inactive older examples and source copies do not define the manuscript
procedure in this comparison.
\end{document}
